%% ==================================================
%% Section: Lieb--Robinson propagation for finite-range Hamiltonians on graphs
%% Sources: Hastings--Koma, arXiv:math-ph/0507008, Appendix A, (A.12)--(A.16);
%% Nachtergaele--Ogata--Sims, arXiv:math-ph/0603064, Section 2; OpenAI, "A
%% two-dimensional area law from a global spectral gap" (September 24, 2026,
%% openai/math@adc7f1241b42e322a6451854ab7e4b4c146bf78a),
%% build/sections/03-quasilocal.tex, lines 14--41 and 52--116 (Lemma 4.1,
%% label lem:quasilocal-lr). The results below are stated for a general
%% finite graph and a general interaction family; the induced square-lattice
%% statement of Lemma 4.1 is a separate entry.
%% ==================================================

\section{Lieb--Robinson propagation on finite graphs}
\label{sec:qc_lieb_robinson}

Let $I$ be a finite set of sites of local dimension $q$, and write
$\mathcal A_X$ for the operators on $(\C^q)^{\otimes I}$ acting on
$X\subseteq I$. All norms are Hilbert-space operator norms. For an operator
$H$ put $\tau_t(A)=e^{itH}Ae^{-itH}$. The dynamics generated by a sum of
local Hermitian terms moves support at a bounded speed, up to exponentially
small tails \cite{Hastings2006Spectral,Nachtergaele2006LiebRobinson}. The
estimate below measures distance by the path metric of an arbitrary simple
graph on $I$, so that holes and separate connected components cost nothing:
sites in different components are at infinite distance, and there the
commutator vanishes exactly.

\begin{definition}[Restricted commutator norm]
    \label{aux:restricted-commutator-norm}
    \lean{QuantumCircuit.heisenbergCommutatorNorm,
        QuantumCircuit.heisenbergEvolution}\leanok
    For operators $H,B$ and a set of sites $X$, define
    \[
        C_B(X,t)=\sup\{\|[\tau_t(A),B]\|: A\in\mathcal A_X,\ \|A\|\leq1\}.
    \]
\end{definition}

\begin{lemma}[Evaluation of the restricted commutator norm]
    \label{aux:restricted-commutator-evaluation}
    \lean{QuantumCircuit.norm_heisenberg_commutator_le_heisenbergCommutatorNorm}\leanok
    \uses{aux:restricted-commutator-norm}
    For $A\in\mathcal A_X$, one has
    $\|[\tau_t(A),B]\|\leq C_B(X,t)\|A\|$.
\end{lemma}
\begin{proof}\leanok
    Apply the defining operator-norm bound to $A$.
\end{proof}

\begin{lemma}[Reversal of time]
    \label{aux:restricted-commutator-negative-time}
    \lean{QuantumCircuit.heisenbergCommutatorNorm_neg_time}\leanok
    \uses{aux:restricted-commutator-norm}
    The restricted commutator norm for $H$ at time $-t$ equals that for
    $-H$ at time $t$.
\end{lemma}
\begin{proof}\leanok
    Both Heisenberg evolutions are $A\mapsto e^{-itH}Ae^{itH}$.
\end{proof}

\begin{lemma}[Continuity in time]
    \label{aux:restricted-commutator-continuity}
    \lean{QuantumCircuit.continuous_heisenbergCommutatorNorm}\leanok
    \uses{aux:restricted-commutator-norm}
    For fixed $H,B,X$, the function $t\mapsto C_B(X,t)$ is continuous.
\end{lemma}
\begin{proof}\leanok
    Matrix exponentiation and multiplication are continuous; the norm of the
    resulting linear map on $\mathcal A_X$ is continuous.
\end{proof}

\begin{lemma}[Initial commutator bound]
    \label{aux:restricted-initial-commutator-bound}
    \lean{QuantumCircuit.heisenbergCommutatorNorm_zero_le}\leanok
    \uses{aux:restricted-commutator-norm}
    For arbitrary $H,B,X$, one has $C_B(X,0)\leq2\|B\|$.
\end{lemma}
\begin{proof}\leanok
    Use $\|AB-BA\|\leq2\|A\|\|B\|$.
\end{proof}

\begin{lemma}[Vanishing initial commutator]
    \label{aux:restricted-initial-commutator-vanishing}
    \lean{QuantumCircuit.heisenbergCommutatorNorm_zero_of_disjoint}\leanok
    \uses{aux:restricted-commutator-norm}
    If $B\in\mathcal A_Y$ and $X\cap Y=\varnothing$, then $C_B(X,0)=0$.
\end{lemma}
\begin{proof}\leanok
    Operators on disjoint tensor factors commute.
\end{proof}

\begin{lemma}[Local commutator recursion]
    \label{aux:restricted-commutator-recursion}
    \lean{QuantumCircuit.heisenbergCommutatorNorm_le_integral}\leanok
    \uses{aux:restricted-commutator-norm}
    Let $H=\sum_j h_j$ be a finite sum, where each $h_j$ is Hermitian and
    belongs to $\mathcal A_{T_j}$. For arbitrary $B,X$ and $t\geq0$,
    \[
        C_B(X,t)\leq C_B(X,0)
        +2\sum_{j:X\cap T_j\neq\varnothing}\|h_j\|
        \int_0^t C_B(T_j,s)\,ds.
    \]
    No bound on the ground energy or positivity assumption on $h_j$ is
    required, and several terms may share a support. This is the recursion of
    Hastings--Koma, Appendix A, (A.12)--(A.14).
\end{lemma}
\begin{proof}\leanok
    Separate the terms disjoint from $X$. Their sum commutes with every operator
    in $\mathcal A_X$. Differentiate in the corresponding interaction picture,
    apply the fundamental theorem of calculus and invariance of operator norm
    under unitary conjugation, and then take the restricted operator norm.
\end{proof}

\begin{lemma}[Weighted integral comparison]
    \label{aux:weighted-volterra-comparison}
    \lean{finite_volterra_le_weight_mul_exp,
        finite_volterra_zero_le_weight_mul_exp_sub_one,
        finite_rectangular_volterra_zero_le_weight_mul_exp_sub_one,
        le_mul_exp_of_le_const_add_integral}\leanok
    Let $f_i\geq0$, for $i$ in a finite set, be continuous on $[0,T]$ and satisfy
    $f_i(t)\leq b_i+\sum_j a_{ij}\int_0^tf_j(s)\,ds$ with $a_{ij}\geq0$.
    If positive weights $w_i$ satisfy $\sum_ja_{ij}w_j\leq vw_i$ and
    $b_i\leq cw_i$, then $f_i(t)\leq cw_ie^{vt}$ for $0\leq t\leq T$. If moreover
    $b_i=0$ and $v>0$, then $f_i(t)\leq cw_i(e^{vt}-1)$.
\end{lemma}
\begin{proof}\leanok
    The quantity $\max_if_i(t)/w_i$ satisfies the scalar inequality
    $g(t)\leq c+v\int_0^tg$, and Gr\"onwall's inequality gives $g(t)\leq ce^{vt}$.
    For the refinement, insert this bound into the integral once more.
\end{proof}

\begin{lemma}[Exact locality on interaction-closed regions]
    \label{lem:qc_lr_closed_region}
    \lean{QuantumCircuit.commute_heisenbergEvolution_of_closed}\leanok
    Let $H=\sum_jh_j$ with $h_j\in\mathcal A_{T_j}$, and let $C$ be a set of sites
    such that each $T_j$ is contained in $C$ or disjoint from $C$. If
    $A\in\mathcal A_C$ and $B\in\mathcal A_D$ with $C\cap D=\varnothing$, then
    $[\tau_t(A),B]=0$ for every $t$.
\end{lemma}
\begin{proof}\leanok
    Write $H=H_C+H'$, with $H_C$ the sum of the terms inside $C$. The two parts
    commute, and $H'$ commutes with $\mathcal A_C$, so
    $\tau_t(A)=e^{itH_C}Ae^{-itH_C}$. Every factor commutes with $B$.
\end{proof}

\paragraph{Finite-range interactions on a graph.}
Let $G$ be a simple graph on $I$ with path metric $d_G$, taking the value
$\infty$ between different connected components. Let $H=\sum_kh_k$ be a
finite sum of Hermitian terms with $h_k\in\mathcal A_{X_k}$, and choose an
anchor $a_k\in X_k$. Assume that every support has $d_G$-diameter at most
$R$ and at most $v_R$ sites, and that the norms of the terms whose support
contains any given site sum to at most $b_0$.

\begin{theorem}[Graph-distance propagation]
    \label{thm:qc_graph_lieb_robinson}
    \lean{QuantumCircuit.norm_heisenberg_commutator_le_graphDistance,
        QuantumCircuit.exists_graph_lieb_robinson}\leanok
    \uses{aux:restricted-commutator-recursion,aux:weighted-volterra-comparison,
        aux:restricted-commutator-evaluation,aux:restricted-commutator-negative-time,
        aux:restricted-initial-commutator-bound,
        aux:restricted-initial-commutator-vanishing,
        aux:restricted-commutator-continuity,thm:qc_graph_lieb_robinson_zero}
    In this setting let $A\in\mathcal A_{X_i}$, let $y$ be a site, and let
    $B\in\mathcal A_{\{y\}}$. For every $\mu\geq0$, every integer
    $0\leq n\leq d_G(a_i,y)$ and every real $t$,
    \[
        \|[\tau_t(A),B]\|\leq 2e^{\mu R}\|A\|\|B\|
        \exp\bigl(2v_Rb_0e^{2\mu R}|t|-\mu n\bigr).
    \]
    Consequently there are positive constants $C,v,c$, depending only on
    $R,v_R,b_0$, such that
    $\|[\tau_t(A),B]\|\leq C\|A\|\|B\|e^{v|t|-cd_G(a_i,y)}$ whenever
    $d_G(a_i,y)<\infty$, and $[\tau_t(A),B]=0$ otherwise. The constants do not
    depend on the number of sites, on $q$, or on the graph.
\end{theorem}
\begin{proof}\leanok
    Let $F_k(t)=C_B(X_k,t)$ for $t\geq0$. The local commutator recursion gives
    $F_k(t)\leq F_k(0)+\sum_ja_{kj}\int_0^tF_j$, with $a_{kj}=2\|h_j\|$ when
    $X_j$ meets $X_k$ and $a_{kj}=0$ otherwise. Put
    $D(x)=\min\{n,d_G(x,y)\}$ and $w_k=e^{-\mu D(a_k)}$. If $X_j$ meets
    $X_k$, then $d_G(a_k,a_j)\leq2R$, so $w_j\leq e^{2\mu R}w_k$; the norms of the
    terms meeting $X_k$ sum to at most $|X_k|b_0\leq v_Rb_0$. Hence
    $\sum_ja_{kj}w_j\leq2v_Rb_0e^{2\mu R}w_k$. The initial value $F_k(0)$
    vanishes unless $y\in X_k$, and then $D(a_k)\leq R$ and
    $F_k(0)\leq2\|B\|\leq2\|B\|e^{\mu R}w_k$. The weighted integral comparison
    gives $F_i(t)\leq2\|B\|e^{\mu R}w_ie^{2v_Rb_0e^{2\mu R}t}$ with
    $w_i=e^{-\mu n}$. Replace $H$ by $-H$ for negative times and multiply by
    $\|A\|$. For the constants take $\mu=1$; the infinite-distance case is
    Lemma~\ref{thm:qc_graph_lieb_robinson_zero}.
\end{proof}

\begin{lemma}[Vanishing between connected components]
    \label{thm:qc_graph_lieb_robinson_zero}
    \lean{QuantumCircuit.heisenberg_commutator_eq_zero_of_edist_eq_top}\leanok
    \uses{lem:qc_lr_closed_region}
    In the setting of Theorem~\ref{thm:qc_graph_lieb_robinson}, if
    $d_G(a_i,y)=\infty$ then $[\tau_t(A),B]=0$ for every $t$.
\end{lemma}
\begin{proof}\leanok
    Let $C$ be the connected component of $a_i$. A support of finite diameter
    lies in one component, so it is contained in $C$ or disjoint from it; also
    $X_i\subseteq C$ and $y\notin C$. Apply
    Lemma~\ref{lem:qc_lr_closed_region}.
\end{proof}

\paragraph{Normalized partial traces.}
For $K\subseteq I$ let
$E_K(B)=q^{-|I\setminus K|}\operatorname{Tr}_{I\setminus K}(B)\otimes\mathbb 1_{I\setminus K}$,
with every tensor factor at its original site.

\begin{definition}[Normalized partial-trace expectation]
    \label{def:qc_site_expectation}
    \lean{QuantumCircuit.siteExpectation, QuantumCircuit.traceOutside,
        QuantumCircuit.siteExpectationLM}\leanok
    The map $E_K$ on the operators of $(\C^q)^{\otimes I}$ defined above.
\end{definition}

\begin{lemma}[Weyl averaging formula]
    \label{lem:qc_site_expectation_average}
    \lean{QuantumCircuit.siteExpectation_eq_average,
        QuantumCircuit.sum_localWeyl_mul_star, QuantumCircuit.outsideWeyl}\leanok
    \uses{def:qc_site_expectation}
    Let $q\geq1$, and for $p=(a,b)$ let $W(p)=X^aZ^b$ be the Weyl operators on
    $\C^q$. Then $E_K(B)$ is the uniform average of $W_pBW_p^\dagger$ over all
    products $W_p=\bigotimes_{x\notin K}W(p_x)$ of Weyl operators outside $K$.
\end{lemma}
\begin{proof}\leanok
    The Weyl operators satisfy
    $\sum_pW(p)_{ii'}\overline{W(p)_{jj'}}=q\,\delta_{ij}\delta_{i'j'}$. The
    corresponding sum over the products factorizes over the sites, and the
    resulting Kronecker deltas reproduce the partial-trace formula.
\end{proof}

\begin{lemma}[Properties of the expectation]
    \label{lem:qc_site_expectation_properties}
    \lean{QuantumCircuit.siteExpectation_mem_supportedOperators,
        QuantumCircuit.siteExpectation_of_mem_supportedOperators,
        QuantumCircuit.siteExpectation_one, QuantumCircuit.norm_siteExpectation_le,
        QuantumCircuit.siteExpectationLM_isKrausCPTP}\leanok
    \uses{lem:qc_site_expectation_average}
    For $q\geq1$, the map $E_K$ takes values in $\mathcal A_K$, fixes every
    operator in $\mathcal A_K$, is unital, completely positive and trace
    preserving, and is a contraction in operator norm.
\end{lemma}
\begin{proof}\leanok
    The values act on $K$ by the partial-trace formula. The averaging formula
    is a Kraus decomposition by rescaled unitaries. Operators acting on $K$
    commute with every $W_p$. A uniform average of unitary conjugations is a
    contraction.
\end{proof}

\begin{lemma}[Localization error through on-site commutators]
    \label{lem:qc_site_expectation_commutator}
    \lean{QuantumCircuit.norm_sub_siteExpectation_le}\leanok
    \uses{lem:qc_site_expectation_average}
    If $\|[B,U]\|\leq\varepsilon_y$ for every site $y\notin K$ and every unitary
    $U\in\mathcal A_{\{y\}}$, then $\|B-E_K(B)\|\leq\sum_{y\notin K}\varepsilon_y$.
\end{lemma}
\begin{proof}\leanok
    By the averaging formula, $B-E_K(B)$ is the average of
    $(BW_p-W_pB)W_p^\dagger$. Write $W_p$ as a product of on-site unitaries and
    use $\|[B,UV]\|\leq\|[B,U]\|+\|[B,V]\|$ for unitaries $U,V$.
\end{proof}

\begin{theorem}[Localization onto graph balls]
    \label{thm:qc_graph_ball_localization}
    \lean{QuantumCircuit.norm_heisenberg_sub_siteExpectation_graphBall_le,
        QuantumCircuit.graphBall, QuantumCircuit.sum_pow_mul_exp_neg_le}\leanok
    \uses{thm:qc_graph_lieb_robinson,thm:qc_graph_lieb_robinson_zero,
        lem:qc_site_expectation_commutator,lem:qc_site_expectation_properties}
    In the setting of Theorem~\ref{thm:qc_graph_lieb_robinson}, with $q\geq1$,
    assume that at most $K_g(d+1)^k$ sites lie at distance exactly $d$ from any
    site. Let $N_l(a)=\{x:d_G(a,x)\leq l\}$. Then
    \[
        \|\tau_t(A)-E_{N_l(a_i)}(\tau_t(A))\|
        \leq\|A\|\min\bigl\{2,\;Ce^{v|t|-l/2}\bigr\},
    \]
    with $v=2v_Rb_0e^{2R}$ and $C=2e^RK_g\,2^kk!\,e^{1/2}/(1-e^{-1/2})$.
\end{theorem}
\begin{proof}\leanok
    The bound $2\|A\|$ follows from contractivity. For the other bound, apply the
    commutator lemma with $\varepsilon_y=0$ when $d_G(a_i,y)=\infty$ and
    $\varepsilon_y=2e^R\|A\|e^{v|t|-d_G(a_i,y)}$ otherwise, using graph-distance
    propagation with $\mu=1$. Group the sites outside the ball by their distance
    $d>l$, and use $(d+1)^k\leq2^kk!\,e^{(d+1)/2}$ and a geometric series.
\end{proof}
